\textbf{Puestos fijos y simultaneidad.} Se proponen dieciséis puestos instalados con dos monitores cada uno: siete para administración, socios, contratos y cobranza; dos para gerencia y comercial; dos para torre/radio; dos para portería; uno para coordinación de escuela en tierra; uno para planificación de varadero; y uno para registro de combustible en una oficina segura fuera de la zona clasificada. La dotación de administración de cinco a siete y la de gerencia/comercial de dos proceden del caso. Los dos puestos de torre separan escucha/maniobra y preparación de legajos; los dos de portería separan control de acceso y atención/retiros. Su simultaneidad constituye un supuesto de peak, no una dotación medida. Los turnos de cuatro a seis personas en torre y cinco a siete en acceso no se convierten en un PC por cada persona. El uso de mantenimiento desde el puesto de coordinación de varadero se valida sin desplazar al operador; si exige simultaneidad adicional, se amplían cantidades y energía antes de compra \parencite[cap.~2, secc.~2.4; cap.~14, secc.~14.2]{caso07}.

La Tabla~\ref{tab:p4-puestos} distingue simultaneidad y criticidad de negocio. Los siete puestos críticos mantienen las funciones de recepción, zarpe, escuela y combustible de RT-03.10. La protección eléctrica comprende los dieciséis PC y sus treinta y dos monitores, también los administrativos y de gerencia: la prioridad funcional no reduce la cobertura del suministro ininterrumpido. El balance incluye todos los puestos en régimen normal y con cada rama retirada, junto con sus dependencias. Esta continuidad se ensaya con la Célula Local y las autoridades habilitantes \parencite[RT-06.07, p.~15; RT-08.04, p.~18]{pucv-transversal}.

\begin{table}[htbp]
\centering\small
\caption{Dimensionamiento propuesto de puestos fijos por función simultánea}\label{tab:p4-puestos}
\begin{tabular}{L{\dimexpr0.46\linewidth-2\tabcolsep\relax}R{\dimexpr0.18\linewidth-2\tabcolsep\relax}R{\dimexpr0.18\linewidth-2\tabcolsep\relax}R{\dimexpr0.18\linewidth-2\tabcolsep\relax}}
\hline
\EncabezadoTabla{Función} & \EncabezadoTabla{PC instalados} & \EncabezadoTabla{Monitores} & \EncabezadoTabla{PC críticos} \\ \hline
\rowcolor{SynSurface}
Administración y cobranza & 7 & 14 & 0 \\ \hline
Gerencia y comercial & 2 & 4 & 0 \\ \hline
\rowcolor{SynSurface}
Torre y radio & 2 & 4 & 2 \\ \hline
Portería y atención & 2 & 4 & 2 \\ \hline
\rowcolor{SynSurface}
Escuela en tierra & 1 & 2 & 1 \\ \hline
Varadero en oficina & 1 & 2 & 1 \\ \hline
\rowcolor{SynSurface}
Combustible en oficina segura & 1 & 2 & 1 \\ \hline
Total instalado & 16 & 32 & 7 \\ \hline
\rowcolor{SynSurface}
Stock compatible adicional & 2 & 4 & 0 \\ \hline
Total a proveer & 18 & 36 & 7 \\ \hline
\end{tabular}
\FuenteTabla{elaboración propia: supuesto de simultaneidad derivado de \parencite[cap.~2, secc.~2.4; cap.~14, secc.~14.2; cap.~15, RT-03.10]{caso07}; reserva por tipo según \parencite[secc.~8.4, pp.~19--20]{pucv-transversal}.}
\end{table}

Las reservas se calculan por tipo compatible y son adicionales al instalado: $\lceil0,10\times16\rceil=2$ PC y $\lceil0,10\times32\rceil=4$ monitores. Cada puesto recibe teclado, mouse, dos cables DisplayPort y conexión Ethernet; se añaden dos teclados, dos mouse, cuatro cables DisplayPort y dos juegos de alimentación PC compatibles de stock. Las dieciocho computadoras, incluidas las dos reservas precargadas, quedan inventariadas y gestionadas. Este parque se distingue de los catorce móviles operacionales y dos reservas móviles, que conservan su distribución. La hipótesis de cincuenta sesiones internas de SD2 A2.7, EST-05, sigue siendo escenario de capacidad: treinta terminales operacionales nuevos, dieciséis PC y catorce móviles, no acreditan cincuenta operadores físicos simultáneos. Las veinte sesiones restantes del ensayo son carga adicional simulada, sin inventar terminales existentes. Los puestos remotos técnicos pertenecen a una población de servicio distinta y no aumentan el parque de negocio del CLIENTE.

\textbf{Referencia de estación y selección.} La referencia provisional para dimensionamiento es Dell OptiPlex Micro Plus 7020 con Intel Core i5-14500T de catorce núcleos y veinte hilos, 32 GB DDR5-4800 en dos módulos de 16 GB, SSD M.2 2230 PCIe Gen4 NVMe Class 35 de 512 GB, gráficos Intel UHD 770, Ethernet de 1 Gb/s, TPM 2.0, UEFI y Secure Boot. Sus tres DisplayPort 1.4a permiten conectar dos pantallas independientes. Se evalúa el OptiPlex Micro 7020 básico, que dispone de HDMI y DisplayPort y también admite monitores duales; Plus aporta conexión homogénea DisplayPort y ampliación de video; su fuente única y el NVMe único impiden autorizar esta referencia como configuración final frente a RT-08.04 y RT-08.02. Se descarta un escritorio virtual dependiente exclusivamente de nube para los siete puestos críticos por la continuidad sin WAN. Son criterios de diseño, no un ensayo comparativo realizado \parencites{p4-dell-cpu,p4-dell-memoria,p4-dell-ssd,p4-dell-pantallas,p4-dell-seguridad,p4-dell-micro-alternativa}.

Cada puesto incorpora dos Dell P2425H IPS de 23,8 pulgadas y 1920$\times$1080, con base regulable y cable DP individual. La referencia internacional es 210-BMGH/H127Y; el teclado cableado es KB216 español QWERTY, 580-ADGS/8444F, y el mouse cableado MS116, 570-AAIS/JCYP0. Estos códigos no acreditan SKU ni disponibilidad en Chile. La variante P2425Ht tiene registro ENERGY STAR ID 3999653; su correspondencia con el monitor de compra debe comprobarse. El certificado de monitor no se extiende al PC ni al puesto completo. La BOM concreta de SSD, memoria, adaptador, cables eléctricos locales, montaje, periféricos de atención y mobiliario requiere confirmación compatible; no se presenta una combinación de opciones del manual como una compra ya aprobada \parencites{p4-dell-monitor,p4-energystar-monitor,p4-dell-teclado,p4-dell-mouse}.

La alternativa de sustitución evaluada es Vecow ECS-4700-2G-1345UE con dos SATA Innodisk de 512 GB en RAID 1, dos EIZO MDF2701W-BK y dos fuentes QUINT independientes. El detalle de BOM, alimentación cruzada, tolerancia, aceptación y cantidades se desarrolla en la evaluación complementaria del T-11. Su incorporación sustituiría las mismas unidades del parque CLIENTE y remoto, sin añadir otro parque. La certificación energética de las variantes y la evaluación ergonómica deben completar la selección antes de adquisición; mientras tanto se conservan las cantidades y cotas de dimensionamiento, sin atribuir cumplimiento integral a la candidata \parencites{p9-vecow-datasheet-2025,p6-puestos-eizo}[RT-08.02/04/08, pp.~18--19]{pucv-transversal}.

El diseño y la recepción del puesto comprobarán mesa, silla, alcance, postura, iluminación y ubicación de las pantallas frente al texto exigido por RT-08.08. No se ha localizado una fuente oficial inequívoca de NCh 2527 ni se aportan cláusulas o evaluación individual, por lo que se mantiene esa brecha. Los ajustes de monitor son evidencia parcial y no una acreditación ergonómica. La guía de ergonomía del ISP es apoyo técnico, sin sustituir la norma requerida. Tampoco se acredita todavía certificación energética vigente de la configuración PC elegida o su SKU local \parencites[RT-08.08, p.~19]{pucv-transversal}{p4-isp-ergonomia}.

\textbf{Gestión y protección de estaciones.} Synaptix propone Microsoft Intune Plan 1 for Devices para las dieciocho computadoras compartidas, con enrolamiento por método elegible, grupos de dispositivos, inventario, configuración y despliegue de aplicaciones. La modalidad sin afinidad de usuario admite políticas dirigidas al dispositivo; excluye acceso condicional, protección de aplicaciones y funciones de gestión por usuario. Se verifican método de enrolamiento, requisitos Entra/Autopilot cuando correspondan y licencias efectivas antes de activar. Los controles de identidad funcional Cognito/local mantienen usuarios individuales y no se dan por implementados con esta compra \parencite{p4-intune-licencias}.

Se propone Sophos Endpoint con XDR para protección, investigación/respuesta y control de periféricos, y Sophos Central Device Encryption adicional para BitLocker sobre volúmenes fijos. Se bloquean almacenamiento extraíble y MTP por defecto, con excepciones aprobadas por identificador y uso; los adaptadores de red no se bloquean sin validar su excepción. Live Response queda restringido a roles autorizados y auditado. GPT, UEFI, TPM y cifrado se comprueban antes de habilitar; la custodia de claves y recuperación sin WAN requieren procedimiento verificable. Se propone TPM-only en los siete puestos críticos para permitir reinicio automático sin intervención TI del CLIENTE, con arranque seguro, cifrado y acceso individual a Windows/aplicación; las consolas técnicas remotas usan TPM+PIN con personal autorizado. La política y la custodia de recuperación se comprueban junto con el riesgo físico y los reinicios operacionales. El cifrado no sustituye al inicio de sesión individual y Device Encryption no cifra medios extraíbles. Intune no aplica una segunda política BitLocker contradictoria. UEM Android, protección Linux y este agente Windows mantienen alcances separados \parencites{p4-sophos-perifericos,p4-sophos-respuesta,p4-sophos-cifrado,p4-sophos-politica-cifrado}.

La operación técnica remota de Synaptix especifica dos consolas simultáneas: una para vigilancia y gestión de incidentes, otra para supervisión y escalamiento. Cada una usa la misma configuración OptiPlex Micro Plus 7020 y dos P2425H; una tercera computadora y dos pantallas compatibles forman reserva del servicio. Se incluyen tres teclados y mouse, seis cables DP y alimentación compatible. Son dos puestos técnicos más un juego de reserva, distintos de los dieciséis de negocio y del stock del CLIENTE. Synaptix los provee dentro del servicio gestionado, fuera de la marina; no se atribuye disponibilidad existente ni dotación 24 por 7 por contar consolas. El puesto de visita se propone en la zona segura de administración, usando un PC de reserva solo durante la intervención autorizada; mientras se use, se registra la disminución del stock y se aporta sustituto si es necesario, sin contabilizarlo simultáneamente como reserva disponible. La sala remota, comunicaciones y dotación se individualizan con el modelo de servicio; RT-06.29 conserva su brecha de emplazamiento.

La gestión cubre veintiuna computadoras: dieciocho del CLIENTE y tres del servicio remoto. Las tres técnicas usan la misma gestión y protección, pero tenant, licencias y custodia de Synaptix separados del CLIENTE; no se activa acceso administrativo con cuentas compartidas. La cobertura técnica no equivale automáticamente a veintiuna licencias Sophos de usuario. Se resolverán la modalidad elegible y el padrón de personas beneficiarias, incluyendo técnicos y rotación estacional; prevalece la EULA sobre el cálculo de consola. No se presenta contrato celebrado ni compatibilidad integral comprobada. RT-08.09 conserva No por configuración, licencias y recuperación todavía incompletas. RT-08.07 se formula sobre cantidades y características de ambos parques, sin atribuir instalación realizada \parencites{p4-sophos-licencias}[RT-08.07/09, p.~19]{pucv-transversal}.

Windows 11 Pro 26H2 es línea inicial de referencia, con disponibilidad general desde el 29 de septiembre de 2026 y fin de actualizaciones el 10 de octubre de 2028. Su adopción exige verificar compatibilidad OEM, agentes y aplicaciones; 25H2 es alternativa temporal soportada hasta el 12 de octubre de 2027. Se planifican migraciones y renovaciones hasta mes 56, sin atribuir ese horizonte a una versión única. Intune aplica anillos de actualización por dispositivo, piloto en reservas, validación y reversión; agentes Sophos y BIOS/controladores usan sus mecanismos verificados. Los reinicios de puestos críticos se escalonan y no retiran ambos puestos torre/portería simultáneamente. La periodicidad ordinaria no amplía los máximos urgentes de siete, quince y treinta días. Las políticas precargadas siguen siendo protección local al perder WAN; nuevas políticas y respuesta remota requieren canal \parencites{p4-windows-versiones,p4-intune-actualizaciones}[RT-11.04, p.~23]{pucv-transversal}.

\textbf{Consumo y condiciones de continuidad.} Dell declara adaptadores PC de 130/180 W de salida, sin acreditar la entrada de la BOM. Para reconciliar todos los estados con el balance de SD4 y T-11 se fijan cotas de adquisición de 220 W AC por PC y 65 W por pantalla, con conversión/cableado incluidos una vez: 350 W por PC y dos monitores. Se reservan otros 20 W por puesto para auxiliares externos alimentados que la BOM final requiera; no son teclado/mouse USB ni una medición. La envolvente es 370 W/puesto, 2,590 kW para los siete prioritarios y 5,920 kW para los dieciséis. Una BOM que exceda su cota exige sustitución o recálculo antes de compra. El stock apagado no integra régimen; las cargas exteriores no se imputan automáticamente al calor/PUE de la sala \parencites{p4-dell-adaptador,p4-dell-monitor}.

Los PC y pantallas Dell utilizados como referencia tienen fuente única: una transferencia A/B protege ante falla de rama pero no crea redundancia electrónica. RT-08.04 conserva su brecha literal, y la contingencia y reemplazo compatibles se comprueban con el servicio. El puesto interior de varadero tampoco sustituye cobertura de sus móviles: se propone un gabinete de borde propio y dos switches y dos AP adicionales, conectados a los core por caminos segregados y con energía A/B. Es una hipótesis de cobertura, no un resultado RF. Se mantienen diez concentradores de medición en pantalán: no se agrega uno a varadero sin una función de medición que lo requiera. Así resultan seis gabinetes de borde, doce switches y doce AP; el stock compatible de switches/AP pasa a dos por tipo. Modelo, mapa de rutas/puertos, RF, fuentes y prestaciones después de falla aún requieren individualización. Las nuevas cargas y dispositivos deben integrar el balance de UPS, bancos, generación y climatización antes de compra y aceptación \parencites[RT-08.04, p.~18]{pucv-transversal}[cap.~15, RT-03.24 y RT-06.01]{caso07}.

